\section*{Summary}
\addcontentsline{toc}{section}{Summary} % adds to toc

The main objectives of this thesis were to provide a self-contained proof of Theorem \ref{theor:main_theorem}, identify sufficient conditions on the kernel function for the theorem to hold, and illustrate how diffusion maps perform on both synthetic and real-world data. 

In Definitions and Preliminaries (see Section \ref{sec:one}), we recalled necessary definitions and results from linear algebra and probability theory. 

In Diffusion Distance and Diffusion Maps (see Section \ref{sec:two}), we first proved the spectral theorem, which states that every real symmetric matrix has real eigenvalues and its eigenvectors form an orthonormal basis of $\R^n$. We then developed the theory of diffusion maps in a self-contained way. We stated explicit sufficient conditions on the kernel (strict positivity, symmetry, and positive semi-definiteness) and proved Theorem \ref{theor:main_theorem}, which states that the diffusion distance is equivalent to the Euclidean distance in the embedded space. The original paper by Nadler et al. \cite{nips:dif_maps} proves the theorem using the Gaussian kernel but does not explicitly state the conditions under which the result holds for general kernels and relies on auxiliary results that are not always stated.

In Applications (see Section \ref{sec:three}), we applied the theory to three synthetic datasets and a real-world dataset. The synthetic examples demonstrated how diffusion maps work at an intuitive level (Subsection \ref{subsubsec:db}); how they can be used for dimensionality reduction, including how they differ from PCA, a classical linear method (Subsection \ref{subsec:toroidal helix}); and, lastly, how they can be used to solve correspondence problems (Subsection \ref{subsec:tori}). We concluded the section by applying diffusion maps to two figures from the Dynamic FAUST dataset. Since the Gaussian kernel is not invariant under non-rigid transformation, the embeddings did not align as cleanly as in the synthetic examples. However, when embedding the figures into $\R^3$, the results were similar. Checking the correspondence numerically confirmed that diffusion maps also perform well with complex real-world data.

There are several directions in which this work could be extended. We identified sufficient conditions on the kernel, but we believe they can be weakened: strict positivity is used to guarantee irreducibility and aperiodicity, positive semi-definiteness to guarantee non-negative eigenvalues. We believe weaker conditions that ensure the same properties would suffice. Identifying exact necessary and sufficient conditions would be the natural next step. Furthermore, the Dynamic FAUST example (Subsection \ref{subsec:FAUST}) raises the question of which kinds of kernels are better suited for solving non-rigid correspondence problems. Finally, providing rigorous mathematical justifications for the use of diffusion maps for dimensionality reduction would be a valuable direction for future work. In particular, formalising the spectral-gap argument, which is one of the key criteria for determining the minimal number of dimensions to keep.